\chapter{平面双足连续行走}
任务分层与接触逆动力学参考PlaCo\cite{placo}及TSID\cite{tsid}，不整体复制其控制栈。


将规划参考、实际摆腿、接触和有限电机力矩放入同一物理系统。比较统一的质心与零力矩点接口，同时记录误差与完成率之间的折中。

\section{力矩驱动的平面双足与八步验证}

\label{stage:s21}

\subsection{关节力矩驱动的平面双足}\label{s21-detail-1}

\subsubsection{模型与控制接口}\label{s21-detail-2}

新建27.2kg的平面双足，包含躯干、左右大腿、小腿和脚掌。基座有前后、竖直和平面俯仰三个自由度且没有执行器；两侧各有髋、膝、踝三个电机。每段腿长0.4m，足底约0.20\(\times\)0.09m。侧向、滚转和偏航受模型约束，躯干及连杆碰撞做了简化。

逆运动学只用于生成初始站姿。之后控制器读取物理状态，建立全身逆动力学QP，求解关节加速度、电机力矩和预测接触扳手。模拟器仅接收六个电机力矩；QP预测的接触力不直接施加到模型，基座位置和关节轨迹也不被逐帧改写。实际脚掌接触由MuJoCo求解。相关代码在 models/biped.xml、scripts/biped.py。

\begin{Verbatim}[breaklines=true,breakanywhere=true,fontsize=\footnotesize]
M(q) qdd + bias(q, qd) = S^T τ + J(q)^T λ + passive
约束：力矩上下限、关节位置和预测速度、支撑足加速度、摩擦与足底力矩界
目标：质心、躯干俯仰、摆脚轨迹，以及小幅姿态正则项
\end{Verbatim}

使用OSQP 1.0.4，固定数值容差和自适应rho更新间隔。名义控制100Hz、物理积分1ms。电机力矩上限为每侧80、80、30N·m；预测关节速度界8rad/s，实际超10rad/s终止。接触约束包含非负法向力、\textbar Fx\textbar{}\(\leq\)\(\mu\)Fz与\textbar My\textbar{}\(\leq\)0.09Fz，另限制下一控制时刻的足端速度。

\stagefigure{figures/ch07/s21_step_sequence_6db95e66.png}{同一动力学运行在0.5、2.7、5.4s的真实相机帧。}{fig:s21-detail-1}

\subsubsection{标准行走与推扰}\label{s21-detail-3}

\FloatBarrier\begingroup\small\setlength{\tabcolsep}{2.5pt}
\begin{longtable}[]{@{}
  >{\raggedright\arraybackslash}p{(\linewidth - 2\tabcolsep) * \real{0.5254}}
  >{\raggedright\arraybackslash}p{(\linewidth - 2\tabcolsep) * \real{0.4746}}@{}}
\caption{标准行走与推扰：设置与结果}\label{tab:s21-1}\tabularnewline
\toprule\noalign{}
\begin{minipage}[b]{\linewidth}\raggedright
标准条件
\end{minipage} & \begin{minipage}[b]{\linewidth}\raggedright
结果
\end{minipage} \\
\midrule\noalign{}
\endfirsthead
\toprule\noalign{}
\begin{minipage}[b]{\linewidth}\raggedright
标准条件
\end{minipage} & \begin{minipage}[b]{\linewidth}\raggedright
结果
\end{minipage} \\
\midrule\noalign{}
\endhead
\bottomrule\noalign{}
\endlastfoot
步数与步长 & 8步 \(\times\) 0.12m \\
步时 / 双支撑 & 0.6s / 0.12s \\
实测落脚 & 8/8 \\
前进量 & 0.937028m \\
最大躯干俯仰 & 0.0863° \\
最大关节速度 & 2.6206rad/s \\
终端足位最大误差 & 0.791mm \\
\end{longtable}\endgroup

落脚检查同时要求：在计划时刻附近，正确侧脚的水平位置误差小于2.5cm、足底高度绝对值小于8mm，且实测法向力超过10N持续至少15ms。还要求完整运行、前进量不低于计划距离减10cm、终端足位误差小于4cm，并且没有触发控制器或姿态保护。仅保持不倒不算完成行走。

\stagefigure{figures/ch07/s21_measured_steps_d93b29e0.png}{实际足端轨迹与地面接触，用于核对摆脚和八次落脚。}{fig:s21-detail-2}

\subsubsection{条件扫描与失败边界}\label{s21-detail-4}

扫描步长0.08/0.12/0.16m、控制率50/100/200Hz、水平推力0/20/40N、摩擦系数0.2/0.4/0.8、力矩比例0.5/0.75/1、不同双支撑时长及0.5ms积分步长。推扰持续0.1s，自2.70s开始。各组只改变对应实验因素，重复配置去重后为51组；包括4组站立检查。

\stagefigure{figures/ch07/s21_biped_constraints_edf641d4.png}{动力学条件对照；全部条件和结果保存在 \texttt{biped\_results.json}。}{fig:s21-detail-3}

51组中通过50组。唯一未通过条件为步长0.12m、100Hz、双支撑0.12s、摩擦0.8、力矩减半并受40N推扰，在4.191s因躯干姿态或高度越界终止。该例前进约1.083m后失稳，前进得更远并不代表任务完成。完整力矩在同样推扰下完成八步。

\stagefigure{figures/ch07/s21_biped_keyframe_1ca131fc.png}{完整力矩与半力矩对照。右侧在终止后保留真实末帧，未继续模拟。}{fig:s21-detail-4}

该阶段只对速度和力矩幅值设置硬约束，没有实现电机力矩变化率或执行器带宽模型。标准运行在0.5s后相邻控制指令的最大力矩变化约51.5N·m，仍可能对真实驱动器不友好。QP耗时是本机离线记录，不代表硬实时保证。

\subsection{独立检查与实现修正}\label{s21-detail-5}

检查器重新读取保存的轨迹，而不是只采用运行器写出的成功标记。反馈与规划层用独立离散动力学公式重算相邻状态；动力学层从关节状态重新计算足位和质心，并检查实际施加的力矩、关节积分和落脚记录。相关代码在 scripts/verify.py、scripts/tests.py。

\FloatBarrier\begingroup\small\setlength{\tabcolsep}{2.5pt}
\begin{longtable}[]{@{}
  >{\raggedright\arraybackslash}p{(\linewidth - 2\tabcolsep) * \real{0.4751}}
  >{\raggedright\arraybackslash}p{(\linewidth - 2\tabcolsep) * \real{0.5249}}@{}}
\caption{独立检查与实现修正：设置与结果}\label{tab:s21-2}\tabularnewline
\toprule\noalign{}
\begin{minipage}[b]{\linewidth}\raggedright
检查项
\end{minipage} & \begin{minipage}[b]{\linewidth}\raggedright
结果
\end{minipage} \\
\midrule\noalign{}
\endfirsthead
\toprule\noalign{}
\begin{minipage}[b]{\linewidth}\raggedright
检查项
\end{minipage} & \begin{minipage}[b]{\linewidth}\raggedright
结果
\end{minipage} \\
\midrule\noalign{}
\endhead
\bottomrule\noalign{}
\endlastfoot
单元与回归检查 & 11项通过 \\
反馈层因果性 & 无未来测量；接受的源时间戳单调；过期终止符合阈值 \\
反馈层动力学最大残差 & 2.89e-15 \\
规划层动力学最大残差 & 8.88e-16 \\
规划层摆脚 & 成功轨迹连续极值不穿地、速度界满足，换脚不瞬移 \\
动力学基座与关节 & 3个非驱动自由度、6个电机；未发现关节状态重置 \\
动力学正运动学最大残差 & 0.00e+00m \\
关节积分最大残差 & 1.11e-16 \\
成功运行QP动力学最大残差 & 2.84e-04 \\
\end{longtable}\endgroup

实现中修正了两个影响结果的问题。第一，MuJoCo积分后部分运动学缓存仍可能对应上一步状态，直接做雅可比差分会放大Jdot误差；现在先同步运动学，再计算控制量，并加入中央差分回归。第二，摆脚在线改期可能产生很小的负高度，现通过连续极值筛查拒绝该候选，而非裁剪显示。最终批次在这些修正之后重新运行，早期调试结果不混入正式统计。

每个正式批次保存配置、源码哈希与源码快照，校验时核对当前文件与快照一致。录屏是同配置的独立重跑，不增加实验数量。源码包保留配置与汇总结果，原始数组留在本地 runs/。


\paragraph{阶段结论}51个唯一动力学条件通过50个；标准八步前进约0.937米，但尚未统一全部反馈时序实验。



% BEGIN SIM_VISUALS s21
\subsection{实际摆腿与失稳对照}
保留完整力矩和半力矩两侧的原始对照画面。右侧在失败后停留于最后保存姿态，不继续补出运动；“仍显示机器人”不等于该条件继续完成了行走。

\motionfigure{figures/ch07/planar_biped_failure_camera_sequence.png}{完整力矩与半力矩的实际相机画面对照；后者终止后冻结末帧，不能将其显示时长当成有效运行时长。}{fig:motion-planar_biped_failure_camera_sequence}
% END SIM_VISUALS s21
\stagecode{s21}

\stagevideo{videos/ch07/planar_biped_torque_comparison.mp4}

\FloatBarrier

\section{质心与零力矩点到关节力矩的统一接口}

\label{stage:s22}

\subsection{反馈对应的问题与该阶段范围}\label{s22-detail-1}

原先的lipm\_balance把LIPM步态分析与MuJoCo站立恢复分别验证，确实不能把它们合起来称作物理行走。之后的lipm\_extension已补入浮动基座、真实摆腿及八次实测落脚，但三个实验层次仍没有完全连通。

\FloatBarrier\begingroup\small\setlength{\tabcolsep}{2.5pt}
\begin{longtable}[]{@{}
  >{\raggedright\arraybackslash}p{(\linewidth - 4\tabcolsep) * \real{0.2364}}
  >{\raggedright\arraybackslash}p{(\linewidth - 4\tabcolsep) * \real{0.3073}}
  >{\raggedright\arraybackslash}p{(\linewidth - 4\tabcolsep) * \real{0.4563}}@{}}
\caption{反馈对应的问题与该阶段范围：设置与结果}\label{tab:s22-1}\tabularnewline
\toprule\noalign{}
\begin{minipage}[b]{\linewidth}\raggedright
反馈中的问题
\end{minipage} & \begin{minipage}[b]{\linewidth}\raggedright
已有基础
\end{minipage} & \begin{minipage}[b]{\linewidth}\raggedright
该阶段新增
\end{minipage} \\
\midrule\noalign{}
\endfirsthead
\toprule\noalign{}
\begin{minipage}[b]{\linewidth}\raggedright
反馈中的问题
\end{minipage} & \begin{minipage}[b]{\linewidth}\raggedright
已有基础
\end{minipage} & \begin{minipage}[b]{\linewidth}\raggedright
该阶段新增
\end{minipage} \\
\midrule\noalign{}
\endhead
\bottomrule\noalign{}
\endlastfoot
只做站立恢复 & 后续版本已有六电机八步行走 & 同模型延长为20步，加入推扰，并保存完整实际相机录屏 \\
摆脚瞬时交换 & 已有带惯量的腿和连续摆脚参考 & 保留实际动力学与落脚检查，新增计划承重与实测接触核对 \\
ZMP与踝力矩割裂 & 已有质心任务的逆动力学QP & 将ZMP接触力矩平衡显式接入同一个QP，记录预测与实际CoP \\
双支撑需细化 & 已有双足接触和负载目标 & C2负载转移、参考ZMP同步、末态回到两脚承重，扫描不同过渡时长 \\
不同频率与延迟 & 原先只在简化LIPM层测试 & 直接在多关节力矩闭环中比较50/100/200Hz与20/40/80ms水平CoM延迟 \\
\end{longtable}\endgroup

所有新增工作仅在仿真中进行。继续使用27.2kg、9自由度的矢状面模型：三个基座自由度没有执行器，两腿各三个电机。侧向、滚转、偏航受约束；没有使用真实IMU、Jetson或硬件代测。旧目录、旧报告和旧结果保持原样。

\subsection{ZMP到六关节力矩的统一接口}\label{s22-detail-2}

\subsubsection{先统一物理量，再讨论增益}\label{s22-detail-3}

ZMP是地面上的位置，单位为米；踝电机力矩单位为N·m，不能直接把两者数值对应，也不应强行让不同任务的增益相同。

新接口为：

\begin{Verbatim}[breaklines=true,breakanywhere=true,fontsize=\footnotesize]
步态参考CoM位置/速度/加速度
        ↓ 同一组水平质心反馈增益
质心加速度目标 a_cmd
        ↓ 恒高、零质心角动量变化近似
ZMP目标 p_cmd
        ↓ 与CoM任务共同进入全身逆动力学QP
六个电机力矩 τ
        ↓ MuJoCo独立求解真实脚地接触
实际接触力、压力中心和机器人运动
\end{Verbatim}

外环加速度只计算一次：\path|a_cmd = a_ff + 80(c_ref-c_obs) + 18(v_ref-v_obs)|。CoM任务与ZMP目标都使用这个结果，避免两个控制层分别调不同的水平反馈增益。躯干俯仰和摆脚的跟踪增益仍按各自任务设置。

若竖直加速度为a\_z，则采用 \path|p_cmd = c_x - c_z*a_x/(g+a_z)|；超出当前支撑足范围时先投影到支撑范围。每个足底的扳手以足底site为参考、按世界坐标表示为{[}Fx,Fz,My{]}，合成地面压力中心为：

\begin{Verbatim}[breaklines=true,breakanywhere=true,fontsize=\footnotesize]
p = Σ(x_i Fz_i - z_i Fx_i - My_i) / ΣFz_i
Σ[(x_i-p_cmd) Fz_i - z_i Fx_i - My_i] ≈ 0
\end{Verbatim}

式中的My是足底接触力矩，不是踝关节电机力矩。关节力矩由完整动力学约束确定：\path|M qdd + bias = S^Tτ + J^Tλ + passive|。浮动基座没有电机，因此QP必须通过支撑脚接触来实现加速度；求出的接触力不会直接施加到模拟器。

\subsubsection{用软目标保留多关节动力学的差异}\label{s22-detail-4}

实际多关节模型并不严格满足恒质心高度和零角动量变化。ZMP平衡因此作为软目标加入，仍保留关节限位、力矩、摩擦、足底力矩和支撑足加速度等约束。最终运动由MuJoCo推进，不按参考直接改写qpos。

新增的``仅CoM任务''对照与``CoM＋ZMP映射''共用该阶段模型、步态、双支撑、增益和接触核对，区别是是否显式优化ZMP力矩平衡。这是同版本消融，不把旧报告中不同配置的成绩直接拼在一起。

\FloatBarrier\begingroup\small\setlength{\tabcolsep}{2.5pt}
\begin{longtable}[]{@{}
  >{\raggedright\arraybackslash}p{(\linewidth - 4\tabcolsep) * \real{0.2641}}
  >{\raggedright\arraybackslash}p{(\linewidth - 4\tabcolsep) * \real{0.3804}}
  >{\raggedright\arraybackslash}p{(\linewidth - 4\tabcolsep) * \real{0.3555}}@{}}
\caption{用软目标保留多关节动力学的差异：设置与结果}\label{tab:s22-2}\tabularnewline
\toprule\noalign{}
\begin{minipage}[b]{\linewidth}\raggedright
控制频率（Hz）
\end{minipage} & \begin{minipage}[b]{\linewidth}\raggedright
实际CoP RMSE：对照\(\rightarrow\)映射（mm）
\end{minipage} & \begin{minipage}[b]{\linewidth}\raggedright
CoM RMSE：对照\(\rightarrow\)映射（mm）
\end{minipage} \\
\midrule\noalign{}
\endfirsthead
\toprule\noalign{}
\begin{minipage}[b]{\linewidth}\raggedright
控制频率（Hz）
\end{minipage} & \begin{minipage}[b]{\linewidth}\raggedright
实际CoP RMSE：对照\(\rightarrow\)映射（mm）
\end{minipage} & \begin{minipage}[b]{\linewidth}\raggedright
CoM RMSE：对照\(\rightarrow\)映射（mm）
\end{minipage} \\
\midrule\noalign{}
\endhead
\bottomrule\noalign{}
\endlastfoot
50 & 23.86 \(\rightarrow\) 13.51 & 0.64 \(\rightarrow\) 1.37 \\
100 & 19.41 \(\rightarrow\) 6.06 & 0.28 \(\rightarrow\) 1.26 \\
200 & 18.52 \(\rightarrow\) 3.61 & 0.15 \(\rightarrow\) 1.21 \\
\end{longtable}\endgroup

\stagefigure{figures/ch07/s22_mapping_f57b4e97.png}{相同无扰动条件下的压力中心与质心误差。}{fig:s22-detail-1}

100Hz时压力中心误差下降约68.8\%，但质心位置误差由0.28mm增至1.26mm。因此这里能够确认的是映射更明确、压力中心目标跟踪更准，不能概括为所有平衡指标都更好。QP预测CoP也单独保存，不用它替代实际接触测量。

实际CoP在每次新力矩施加前采样，描述当时的真实接触；QP预测对应该次新指令。二者的时间与物理含义不同。压力中心误差统计从0.5s开始；失败运行仅有终止前数据，比较时同时报告是否完成，避免把提前终止的较短轨迹当成更好表现。

相关代码在 \path|scripts/biped.py| 的WBC.solve和 \path|scripts/feedback.py| 的contact\_wrenches、resultant\_cop。


\paragraph{阶段结论}统一水平加速度目标改善实际压力中心跟踪，同时增加质心误差，不是所有指标一起提高。

\stagecode{s22}

\FloatBarrier

\section{连续摆腿、双支撑与二十步行走}

\label{stage:s23}

\subsection{摆腿与双支撑的物理执行}\label{s23-detail-1}

\subsubsection{摆腿不是动画或位置强制驱动}\label{s23-detail-2}

沿用每段0.4m的髋膝踝模型、足底真实碰撞和有限电机力矩。逆运动学只在初始化生成站姿；之后摆脚采用五次多项式水平轨迹与C2离地曲线，目标通过QP转为力矩。名义步长0.12m、步时0.6s、离地高度6cm；物理积分1ms。

每次落脚须满足正确侧脚接近目标位置、足底高度在允许范围内，且实际法向力超过10N持续至少15ms。完整成功还要求运行至结束、完成全部落脚、满足前进量和终端足位误差约束。

两组20步条件均在14s内完成20/20次实测落脚。无扰动与40N推扰下前进量均约2.323m；这是平面模型中的连续步行，不是三维人形走路，也不是只用``不倒''作为成功标准。

\stagefigure{figures/ch07/s23_twenty_steps_11830081.png}{连续20步的实际足位、离地高度和地面法向力。}{fig:s23-detail-1}

\subsubsection{双支撑负载与ZMP共同过渡}\label{s23-detail-3}

使用五次平滑函数 \texttt{s(u)=10u³-15u⁴+6u⁵} 转移计划负载，端点一、二阶导数均为0。参考ZMP按相同负载比例从旧支撑脚中心移到新支撑脚中心；最后一步后渐变到两脚各一半负载，而非一直单脚承重。负载是QP目标，实际负载由接触求解器产生，两者都保留在记录中。

\FloatBarrier\begingroup\small\setlength{\tabcolsep}{2.5pt}
\begin{longtable}[]{@{}
  >{\raggedright\arraybackslash}p{(\linewidth - 6\tabcolsep) * \real{0.2399}}
  >{\raggedright\arraybackslash}p{(\linewidth - 6\tabcolsep) * \real{0.2399}}
  >{\raggedright\arraybackslash}p{(\linewidth - 6\tabcolsep) * \real{0.2484}}
  >{\raggedright\arraybackslash}p{(\linewidth - 6\tabcolsep) * \real{0.2719}}@{}}
\caption{双支撑负载与ZMP共同过渡：设置与结果}\label{tab:s23-1}\tabularnewline
\toprule\noalign{}
\begin{minipage}[b]{\linewidth}\raggedright
双支撑时长（s）
\end{minipage} & \begin{minipage}[b]{\linewidth}\raggedright
40N条件实测落脚
\end{minipage} & \begin{minipage}[b]{\linewidth}\raggedright
最大相邻力矩变化（N·m）
\end{minipage} & \begin{minipage}[b]{\linewidth}\raggedright
双支撑左足比例平均绝对误差（百分点）
\end{minipage} \\
\midrule\noalign{}
\endfirsthead
\toprule\noalign{}
\begin{minipage}[b]{\linewidth}\raggedright
双支撑时长（s）
\end{minipage} & \begin{minipage}[b]{\linewidth}\raggedright
40N条件实测落脚
\end{minipage} & \begin{minipage}[b]{\linewidth}\raggedright
最大相邻力矩变化（N·m）
\end{minipage} & \begin{minipage}[b]{\linewidth}\raggedright
双支撑左足比例平均绝对误差（百分点）
\end{minipage} \\
\midrule\noalign{}
\endhead
\bottomrule\noalign{}
\endlastfoot
0.0 & 8/8 & 42.92 & - \\
0.06 & 8/8 & 9.99 & 14.25 \\
0.12 & 8/8 & 5.96 & 7.00 \\
0.2 & 8/8 & 5.75 & 4.95 \\
\end{longtable}\endgroup

\stagefigure{figures/ch07/s23_double_support_7b86abbd.png}{相同支撑转换附近的计划与实际左足负载比例。}{fig:s23-detail-2}

本组中，无双支撑也能走完八步，但相邻力矩跳变达到42.92N·m；0.12s双支撑下降至5.96N·m。更长过渡使本组负载跟踪更平滑，但会缩短摆腿可用时间，不能据这几组就断言越长越好。零双支撑没有有效的持续过渡区间，不报告其负载比例平均误差。

\subsubsection{接触核对和力矩变化率}\label{s23-detail-4}

支撑集合仍有名义时间表，但计划需要承重的脚必须实际产生法向力；连续超过40ms仍不足5N则终止，不继续假设它在支撑。计划负载低于5\%的卸载脚不触发这一检查。这是对支撑假设的核对，还不是基于触地事件在线重排步时。

另外在QP中约束相邻控制力矩：\path!|τ_k-τ_(k-1)| ≤ rate/hz!。0.4s启动承重后启用，测试600和1500N·m/s，在50/100/200Hz及0/40N推扰的12组条件全部通过，保存的实际指令也逐条核对满足界限。此处限制控制指令斜率，没有建立真实电机电流环或带宽模型。

\stagefigure{figures/ch07/s23_slew_16fc2548.png}{不同控制频率下的实际指令变化率。}{fig:s23-detail-3}

相关代码在 \path|scripts/biped.py| 的Plan.loads、Plan.feet、WBC.solve及simulate。


\paragraph{阶段结论}两组二十步均完成实际落脚并前进约2.323米；名义支撑与实测接触有核对，尚无触地事件重排步时。



% BEGIN SIM_VISUALS s23
\subsection{二十步实际行走的连续画面}
下图来自已完成二十次实测落脚的保存录屏，包含施加40牛、持续0.1秒推力的回合。相机跟随机器人，画面中的地面位置不能直接用来估计前进量；距离和落脚数以前文保存状态与接触记录为准。

\motionfigure{figures/ch07/planar_twenty_steps_full_frame.png}{二十步行走实际相机帧与同期质心、接触记录；画面显示的名义经过时段不代替实测落脚判据。}{fig:motion-planar_twenty_steps_full_frame}

\motionfigure{figures/ch07/planar_twenty_steps_camera_sequence_1.png}{二十步实际动力学录屏的前段相机画面；为保留机器人和脚掌而裁去文字栏。}{fig:motion-planar_twenty_steps_camera_sequence_1}

\motionfigure{figures/ch07/planar_twenty_steps_camera_sequence_2.png}{二十步实际动力学录屏的后段相机画面；为保留机器人和脚掌而裁去文字栏。}{fig:motion-planar_twenty_steps_camera_sequence_2}
% END SIM_VISUALS s23
\stagecode{s23}

\stagevideo{videos/ch07/planar_biped_twenty_steps.mp4}

\FloatBarrier

\section{频率、延迟与动力学独立核验}

\label{stage:s24}

\subsection{频率、推扰和延迟对照}\label{s24-detail-1}

\subsubsection{相同物理模型中的频率与双向推扰}\label{s24-detail-2}

50/100/200Hz控制、0/+40/-40N水平推扰、两种映射，共18个条件；推扰发生在2.70--2.80s。9个仅CoM对照全部完成，9个显式映射完成8个；未通过项为50Hz后向推扰，因计划支撑没有得到实际接触确认而终止。更准的CoP目标跟踪并没有换来更高的条件通过数。

100Hz的推扰恢复如下。恢复定义为推扰结束后，CoM位置误差小于1cm、速度误差小于0.05m/s、躯干俯仰小于0.03rad，并持续至少0.3s。它是相对运动参考的跟踪恢复，不是静止站立恢复；采样分辨率为10ms。若从推扰结束就持续处于阈值内，记0。

\FloatBarrier\begingroup\small\setlength{\tabcolsep}{2.5pt}
\begin{longtable}[]{@{}
  >{\raggedright\arraybackslash}p{(\linewidth - 4\tabcolsep) * \real{0.2550}}
  >{\raggedright\arraybackslash}p{(\linewidth - 4\tabcolsep) * \real{0.3954}}
  >{\raggedright\arraybackslash}p{(\linewidth - 4\tabcolsep) * \real{0.3496}}@{}}
\caption{相同物理模型中的频率与双向推扰：设置与结果}\label{tab:s24-1}\tabularnewline
\toprule\noalign{}
\begin{minipage}[b]{\linewidth}\raggedright
推力（N）
\end{minipage} & \begin{minipage}[b]{\linewidth}\raggedright
推扰后CoM最大误差：对照\(\rightarrow\)映射（mm）
\end{minipage} & \begin{minipage}[b]{\linewidth}\raggedright
恢复时间：对照\(\rightarrow\)映射（s）
\end{minipage} \\
\midrule\noalign{}
\endfirsthead
\toprule\noalign{}
\begin{minipage}[b]{\linewidth}\raggedright
推力（N）
\end{minipage} & \begin{minipage}[b]{\linewidth}\raggedright
推扰后CoM最大误差：对照\(\rightarrow\)映射（mm）
\end{minipage} & \begin{minipage}[b]{\linewidth}\raggedright
恢复时间：对照\(\rightarrow\)映射（s）
\end{minipage} \\
\midrule\noalign{}
\endhead
\bottomrule\noalign{}
\endlastfoot
40 & 6.96 \(\rightarrow\) 5.72 & 0.02 \(\rightarrow\) 0.01 \\
-40 & 10.58 \(\rightarrow\) 14.11 & 0.15 \(\rightarrow\) 0.37 \\
\end{longtable}\endgroup

\stagefigure{figures/ch07/s24_cop_com_c33be67c.png}{40N前向推扰下的实际CoP与CoM误差，橙色为施力区间。}{fig:s24-detail-1}

前向推扰恢复接近，后向推扰下新映射恢复更慢。这提示在当前权重下，压力中心任务与质心、摆腿之间仍有折中；该阶段不再通过不断调权重把失败样例``调没''。

\subsubsection{延迟只施加到明确的一条反馈链}\label{s24-detail-3}

传感器以100Hz生成带时间戳的水平CoM位置和速度，分别加入1mm与0.01m/s高斯噪声；经固定20/40/80ms延迟到达外环。控制器不能读取尚未送达的测量，比较保持过去测量与常速度预测。每个组合使用3个固定噪声种子，两个方法配对使用相同测量噪声。初始化只提供一帧带噪的预启动观测。

关节编码器、用于QP的模型状态、竖直CoM与接触测量仍理想，只有水平CoM外环延迟。不能把这解释成整个机器人在所有传感器均延迟的情况下完成行走。

\FloatBarrier\begingroup\small\setlength{\tabcolsep}{2.5pt}
\begin{longtable}[]{@{}
  >{\raggedright\arraybackslash}p{(\linewidth - 8\tabcolsep) * \real{0.1828}}
  >{\raggedright\arraybackslash}p{(\linewidth - 8\tabcolsep) * \real{0.1828}}
  >{\raggedright\arraybackslash}p{(\linewidth - 8\tabcolsep) * \real{0.1828}}
  >{\raggedright\arraybackslash}p{(\linewidth - 8\tabcolsep) * \real{0.1828}}
  >{\raggedright\arraybackslash}p{(\linewidth - 8\tabcolsep) * \real{0.2687}}@{}}
\caption{延迟只施加到明确的一条反馈链：设置与结果}\label{tab:s24-2}\tabularnewline
\toprule\noalign{}
\begin{minipage}[b]{\linewidth}\raggedright
固定延迟（ms）
\end{minipage} & \begin{minipage}[b]{\linewidth}\raggedright
保持测量：完成
\end{minipage} & \begin{minipage}[b]{\linewidth}\raggedright
常速度预测：完成
\end{minipage} & \begin{minipage}[b]{\linewidth}\raggedright
两方法共同完成对数
\end{minipage} & \begin{minipage}[b]{\linewidth}\raggedright
共同完成时CoM RMSE：保持\(\rightarrow\)预测（mm）
\end{minipage} \\
\midrule\noalign{}
\endfirsthead
\toprule\noalign{}
\begin{minipage}[b]{\linewidth}\raggedright
固定延迟（ms）
\end{minipage} & \begin{minipage}[b]{\linewidth}\raggedright
保持测量：完成
\end{minipage} & \begin{minipage}[b]{\linewidth}\raggedright
常速度预测：完成
\end{minipage} & \begin{minipage}[b]{\linewidth}\raggedright
两方法共同完成对数
\end{minipage} & \begin{minipage}[b]{\linewidth}\raggedright
共同完成时CoM RMSE：保持\(\rightarrow\)预测（mm）
\end{minipage} \\
\midrule\noalign{}
\endhead
\bottomrule\noalign{}
\endlastfoot
20 & 9/9 & 9/9 & 9 & 5.19 \(\rightarrow\) 1.88 \\
40 & 9/9 & 8/9 & 8 & 9.71 \(\rightarrow\) 2.64 \\
80 & 2/9 & 3/9 & 2 & 37.13 \(\rightarrow\) 9.21 \\
\end{longtable}\endgroup

\stagefigure{figures/ch07/s24_delay_3b2aa9af.png}{延迟条件下的完成数，每格3个噪声种子。}{fig:s24-detail-2}

20ms下，两方法均9/9，常速度预测降低了共同完成样本的CoM误差。40ms下，保持测量9/9，预测8/9；预测平均跟踪更准，并不意味着更可靠。80ms下分别只有2/9和3/9通过，这种简单外推不能覆盖较长延迟和接触切换。误差表只在共同完成的配对上比较，不混合提前终止的短轨迹。

全批有9个失败条件在2.70s推扰前已终止，属于延迟、噪声或接触切换自身的问题，不能统称为``抗40N推力失败''。

相关代码在 \path|scripts/feedback.py| 的CoMChannel，以及 \path|scripts/suite.py| 的固定条件矩阵。

\subsection{运行核验与失败边界}\label{s24-detail-4}

\subsubsection{条件计数}\label{s24-detail-5}

\FloatBarrier\begingroup\small\setlength{\tabcolsep}{2.5pt}
\begin{longtable}[]{@{}
  >{\raggedright\arraybackslash}p{(\linewidth - 2\tabcolsep) * \real{0.5013}}
  >{\raggedright\arraybackslash}p{(\linewidth - 2\tabcolsep) * \real{0.4987}}@{}}
\caption{条件计数：设置与结果}\label{tab:s24-3}\tabularnewline
\toprule\noalign{}
\begin{minipage}[b]{\linewidth}\raggedright
条件组
\end{minipage} & \begin{minipage}[b]{\linewidth}\raggedright
组内完成/条件数
\end{minipage} \\
\midrule\noalign{}
\endfirsthead
\toprule\noalign{}
\begin{minipage}[b]{\linewidth}\raggedright
条件组
\end{minipage} & \begin{minipage}[b]{\linewidth}\raggedright
组内完成/条件数
\end{minipage} \\
\midrule\noalign{}
\endhead
\bottomrule\noalign{}
\endlastfoot
mapping & 17/18 \\
feedback & 40/54 \\
support & 8/8 \\
slew & 12/12 \\
integration & 2/2 \\
long\_walk & 2/2 \\
\end{longtable}\endgroup

组间共享的相同配置去重后共94个，不能把上表各行直接相加。无噪声条件没有用不同seed重复增加样本量；只有延迟噪声条件使用3个种子。条件是主动选取的实验网格，不是抽样得到的实际部署成功率。

79个完成，15个未完成：12个因计划承重未被实际接触确认而终止，1个力矩数值约束未通过，1个QP不可行，1个加速度或力矩上下界冲突。全部保留，未只挑成功条件。

开发中曾出现上下界冲突导致OSQP初始化抛出异常、整批中断的问题。现先检查约束界，并将求解器异常记录为该条件的终止结果；正式94组在修复后重新运行，旧开发日志另存，不合并计数。

\subsubsection{独立轨迹核验}\label{s24-detail-6}

重新读取94条运行的qpos、qvel、实际控制指令和接触记录，而不是只采用运行器的success标记。检查包括源码快照、关节积分、足底正运动学、实际接触扳手的坐标/符号、广义接触力、压力中心重算，以及全部落脚判据。

\FloatBarrier\begingroup\small\setlength{\tabcolsep}{2.5pt}
\begin{longtable}[]{@{}
  >{\raggedright\arraybackslash}p{(\linewidth - 2\tabcolsep) * \real{0.6056}}
  >{\raggedright\arraybackslash}p{(\linewidth - 2\tabcolsep) * \real{0.3944}}@{}}
\caption{独立轨迹核验：设置与结果}\label{tab:s24-4}\tabularnewline
\toprule\noalign{}
\begin{minipage}[b]{\linewidth}\raggedright
核验项目
\end{minipage} & \begin{minipage}[b]{\linewidth}\raggedright
最大残差或结果
\end{minipage} \\
\midrule\noalign{}
\endfirsthead
\toprule\noalign{}
\begin{minipage}[b]{\linewidth}\raggedright
核验项目
\end{minipage} & \begin{minipage}[b]{\linewidth}\raggedright
最大残差或结果
\end{minipage} \\
\midrule\noalign{}
\endhead
\bottomrule\noalign{}
\endlastfoot
运行记录与源码快照核对 & 94组 \\
全部记录中的实测落脚 & 690 \\
关节积分残差 & 2.22e-16 \\
足底正运动学残差（m） & 0.00e+00 \\
接触虚功/广义约束力残差 & 1.71e-13 \\
重新求解的接触扳手残差 & 8.38e-12 \\
地面CoP重算残差（m） & 4.18e-15 \\
加速度到原始ZMP映射残差 & 0.00e+00 \\
接口与物理量单元检查 & 7项通过 \\
\end{longtable}\endgroup

接触扳手与运动学重算采用每条运行均匀抽取的12/8个时刻；关节积分、力矩变化率和落脚检查覆盖完整保存轨迹。误差的物理单位随检查项不同，不把它们相互比较。20步视频共280帧、14.0s，逐帧解码无缺帧，录制运行与正式推扰条件的完整轨迹逐元素一致；不增加独立测试数量。

相关代码在 \path|scripts/tests.py|、\path|scripts/verify.py|、\path|scripts/record.py|。


\paragraph{阶段结论}94个去重条件79个通过；80毫秒延迟多数失败，常速度预测不能解决全部接触切换问题。

\stagecode{s24}

\FloatBarrier
